\documentclass[12pt, a4paper]{article}
 % article, report, book, beamer
% tikz
\usepackage{tikz}
\usetikzlibrary{positioning, arrows.meta, shapes}
% colors
\usepackage{xcolor}
\usepackage{titlesec}
% header
\usepackage{fancyhdr}
\usepackage[utf8]{inputenc}
% Basic
\usepackage{amsfonts}
\usepackage{amssymb}
\usepackage{amsmath}
\usepackage{forest}
\usepackage{array}
\renewcommand{\arraystretch}{1.5}
\title{Exercise 01}
\author{LoveKapibarasan}
\date{\today}                   % or specific date

\begin{document}
\maketitle                      % Create title
% Content here
\pagestyle{fancy}
\fancyhf{} 
\fancyhead[L]{\leftmark}     % section name
\fancyhead[R]{\thepage}      % page number
% colors
\titleformat{\section}
  {\color{red}\normalfont\Large\bfseries}  % Blue, Bold
  {\thesection}{1em}{}
\titleformat{\subsection}
  {\color{red}\normalfont\Large\bfseries}  % Blue, Bold
  {\thesection}{1em}{}
%=====
\section { Chebyshev’s inequality and the law of large numbers}

\subsection*{T1}
\textbf{Chebyshev's Inequality}

$P(|X - \mu| < k\sigma) \geq 1 - \frac{1}{k^2}$ for $k > 0$

\textbf{1.}
\begin{align*}
P(|X - 25| \geq 10) &\leq \frac{1}{25} \\
P(15 < X < 35) &\geq 0.96 \\
P(X \leq 35) &\geq 0.96
\end{align*}

\textbf{2.}
\begin{align*}
P(|X - 25| \geq 5) &\leq \frac{1}{6.25} = 0.16 \\
P(20 < X < 30) &\geq 0.84 \\
\therefore P(X \geq 20) &\geq 0.84
\end{align*}

\subsection*{T2}
\begin{align*}
P(40 - b \leq X \leq 40 + b) &\geq 0.95 \\
P(|X - 40| \leq 5k) &\geq 1 - \frac{1}{k^2} \geq 0.95
\end{align*}

Therefore,
\begin{align*}
k &\geq 2\sqrt{5} \\
b &\geq 10\sqrt{5}
\end{align*}

\subsection*{T3}
\begin{align*}
P(75 \leq X \leq 85) &= P(|X - 80| \leq 5 = k\sigma) \geq 0.9 \geq 1 - \frac{1}{k^2} \\
1 - \frac{\sigma^2}{25} &\geq 0.9
\end{align*}

Therefore, $\boxed{\sigma \leq \sqrt{2.5}}$

\subsection*{Ex1}
$\mu \pm k \cdot \sigma$ rule: 68\% ($k = 1$), 95\% ($k = 2$), 99.7\% ($k = 3$) for normal distribution.

\textbf{1.} $U(-1, 1)$

\begin{align*}
\mu &= 0, \quad \sigma^2 = \frac{1}{3} \\
P(\mu - k\sigma \leq X \leq \mu + k\sigma) &= 
\begin{cases} 
k\sqrt{3}, & k\sigma < 1 \\ 
1, & k\sigma \geq 1
\end{cases}
\end{align*}
(Because the range $[-k\sigma, k\sigma]$ extends beyond $[-1, 1]$ when $k\sigma \geq 1$)

\textbf{2.} $X \sim \text{Geometric}(1/2)$

\begin{align*}
\mu &= \frac{1}{p} = 2, \quad \sigma^2 = \frac{1-p}{p^2} = 2
\end{align*}

\subsubsection*{$k = 1$}
\begin{align*}
\mu \pm \sigma &= [0.586, 3.414]
\end{align*}

Therefore,
\begin{align*}
P(X \in \{1, 2, 3\}) &= \frac{1}{2} + \frac{1}{4} + \frac{1}{8} = \frac{7}{8}
\end{align*}

\subsection*{Ex2}
$M_i = c + U_i$, where $E[U_i] = 0$, $\text{Var}(U_i) = 3$, and measurements are independent ($U_i$ is error).

Thus: $E[M_i] = c$, $\text{Var}(M_i) = 3$

Let the sample mean be: $\overline{M}_n = \frac{1}{n} \sum_{i=1}^{n} M_i$

Then: $E[\overline{M}_n] = c$, $\text{Var}(\overline{M}_n) = \frac{\text{Var}(M_i)}{n} = \frac{3}{n}$ based on CLT.

For a 90\% confidence interval:
\begin{align*}
P\left(|\overline{M}_n - c| < \frac{1}{2}\right) &\geq 1 - \frac{\text{Var}(\overline{M}_n)}{(1/2)^2} \geq 0.9 \\
1 - \frac{3/n}{1/4} &\geq 0.9 \\
n &\geq 120
\end{align*}

At least 120 measurements are required to be 90\% confident by Chebyshev's inequality that the sample mean is within 0.5 K of the true melting point.%=====
\end{document}
